% --- IMPORTS ---
\documentclass[leqno]{article}
\usepackage{graphicx}
\usepackage{dsfont}
\usepackage{mathtools}
\usepackage{amssymb}
\usepackage{amsmath}
\usepackage{amsthm}
\usepackage{float}
\usepackage{ragged2e}
\usepackage{tensor}
\usepackage{fancyhdr}
\usepackage{empheq}
\usepackage[most]{tcolorbox}
\usepackage{enumitem}
\usepackage{dirtytalk}
\usepackage{marginnote}
\usepackage{microtype}
\usepackage[
  a4paper,
  left=0.8in,
  right=1in,
  top=1in,
  bottom=0.5in,
  headheight=14pt,
  headsep=0.4in
]{geometry}
\setlength{\parindent}{0cm}
% --- TikZ Packages ---
\usepackage{pgfplots}
\pgfplotsset{compat=1.18}
\usepgfplotslibrary{fillbetween, polar}
\usetikzlibrary{calc, positioning}
\usetikzlibrary{angles, quotes}
\usetikzlibrary{arrows.meta}
\usetikzlibrary{decorations.pathreplacing, decorations.markings}
\usetikzlibrary{patterns}
\usetikzlibrary{intersections}
% --- URL Packages ---
\usepackage[colorlinks=true,linkcolor=red,citecolor=magenta,urlcolor=black]{hyperref}%
\urlstyle{same}
% --- END OF IMPORTS ---

% --- CUSTOM COMMANDS ---
\RenewDocumentCommand{\marks}{o m}{%
  \IfValueTF{#1}%
    {\marginnote{\raggedright\textbf{[#2]}}[#1]}%
    {\marginnote{\raggedright\textbf{[#2]}}}%
}
\newcommand{\vspacerule}[2]{%
  \par\vspace*{#1}%
  \noindent\hrulefill\par
  \vspace*{#2}%
}
\definecolor{scarlet}{RGB}{205, 40, 75}
\newcommand{\questionitem}{%
  \item\phantomsection
  \addcontentsline{toc}{section}{Question \number\value{enumi}}%
}
\renewcommand{\contentsname}{Questions}
% --- END OF CUSTOM COMMANDS ---

% --- HEADER CONFIG ---
\pagestyle{fancy}
\fancyhead{}
\fancyfoot{}
\renewcommand{\headrulewidth}{0.9pt}
\lhead{\textit{Solutions} - Proof by Induction - Divisibility}
\chead{}
\rhead{\href{https://danieljsmith.org}{danieljsmith.org}}
% --- END OF HEADER CONFIG ---

\begin{document}
\vspace*{20pt}
\tableofcontents
\newpage
\begin{enumerate}[
  label=\textbf{\arabic*.},
  leftmargin=*,
  topsep=0pt,
  partopsep=0pt,
  parsep=0pt
]

% ---- Question 1 ----
\questionitem

Prove by induction that \[n^5-n\] is divisible by $5$ for all positive integers $n$ \marks{5}

\vspace*{10pt}

\begin{tcolorbox}[colback=white, colframe=scarlet, title={\textbf{Solution}}, breakable, grow to left by=6mm, grow to right by=10mm]
We prove by induction that \(n^5-n\) is divisible by \(5\) for all positive integers \(n\).

\textbf{Base case:} \(n=1\)

\[
1^5-1=1-1=0
\]

Since \(0\) is divisible by \(5\), the statement is true when \(n=1\).

\textbf{Inductive hypothesis:}

Assume that for some positive integer \(k\), the statement is true for \(n=k\).  
So assume \(k^5-k\) is divisible by \(5\). Equivalently, there exists an integer \(m\) such that

\[
k^5-k=5m
\]

\textbf{Inductive step:}

We must show that \((k+1)^5-(k+1)\) is also divisible by \(5\).

Let \(f(n)=n^5-n\). Then
\[
f(k+1)=f(k)+\bigl(f(k+1)-f(k)\bigr)
\]

Now calculate the difference:
\begin{align*}
f(k+1)-f(k)
&=\bigl((k+1)^5-(k+1)\bigr)-\bigl(k^5-k\bigr)\\
&=(k+1)^5-k^5-1\\
&=\left(k^5+5k^4+10k^3+10k^2+5k+1\right)-k^5-1\\
&=5k^4+10k^3+10k^2+5k\\
&=5(k^4+2k^3+2k^2+k)
\end{align*}

So
\begin{align*}
(k+1)^5-(k+1)
&=(k^5-k)+5(k^4+2k^3+2k^2+k)
\end{align*}

By the inductive hypothesis, \(k^5-k\) is divisible by \(5\).  
Also, \(5(k^4+2k^3+2k^2+k)\) is clearly divisible by \(5\).

Therefore \((k+1)^5-(k+1)\) is divisible by \(5\).

\textbf{Conclusion:}

The statement is true for \(n=1\), and if it is true for \(n=k\), then it is true for \(n=k+1\). Hence, by mathematical induction,

\[
n^5-n \text{ is divisible by } 5 \text{ for all positive integers } n
\]
\end{tcolorbox}


\newpage

% ---- Question 2 ----
\questionitem

Prove by induction that \[9^n-8n-1\] is divisible by $64$ for all positive integers $n\in\mathbb{Z}^+$ \marks{7}

\vspace*{10pt}

\begin{tcolorbox}[colback=white, colframe=scarlet, title={\textbf{Solution}}, breakable, grow to left by=6mm, grow to right by=10mm]
Let
\[
P(n):\quad 9^n-8n-1 \text{ is divisible by } 64
\]

We prove \(P(n)\) is true for all positive integers \(n\) by induction.

\textbf{Base case: \(n=1\)}

\[
9^1-8(1)-1=9-8-1=0
\]

Since \(0\) is divisible by \(64\), \(P(1)\) is true.

\textbf{Inductive hypothesis}

Assume that \(P(k)\) is true for some positive integer \(k\).

So we assume
\[
9^k-8k-1 \text{ is divisible by } 64
\]

Hence there exists an integer \(m\) such that
\[
9^k-8k-1=64m
\]

\textbf{Inductive step}

We now consider
\[
9^{k+1}-8(k+1)-1
\]

\begin{align*}
9^{k+1}-8(k+1)-1
&= 9\cdot 9^k-8k-8-1 \\
&= 9\cdot 9^k-8k-9 \\
&= 9(9^k-8k-1)+64k
\end{align*}

Now use the inductive hypothesis \(9^k-8k-1=64m\):

\begin{align*}
9^{k+1}-8(k+1)-1
&= 9(64m)+64k \\
&= 64(9m+k)
\end{align*}

Since \(9m+k\) is an integer, \(9^{k+1}-8(k+1)-1\) is divisible by \(64\).

Therefore \(P(k+1)\) is true.

\textbf{Conclusion}

Since \(P(1)\) is true, and \(P(k)\Rightarrow P(k+1)\), it follows by induction that
\[
9^n-8n-1 \text{ is divisible by } 64 \text{ for all positive integers } n
\]
\end{tcolorbox}


\newpage

% ---- Question 3 ----
\questionitem

Prove by induction that, for $n \in \mathbb{Z}^+$,
\[13^n+3\cdot 4^n+11\] is divisible by $12$ \marks{7}

\vspace*{10pt}

\begin{tcolorbox}[colback=white, colframe=scarlet, title={\textbf{Solution}}, breakable, grow to left by=6mm, grow to right by=10mm]
\begin{enumerate}[label=\textbf{(\alph*)}]
  \item Let
  \[
  P(n): \text{``}13^n+3\cdot 4^n+11 \text{ is divisible by }12\text{''}
  \]

  We prove \(P(n)\) is true for all \(n\in\mathbb{Z}^+\) by induction.

  \textbf{Base case: \(n=1\).}

  \[
  13^1+3\cdot 4^1+11=13+12+11=36=12\cdot 3
  \]

  So \(13^1+3\cdot 4^1+11\) is divisible by \(12\). Hence \(P(1)\) is true.

  \textbf{Inductive hypothesis.}

  Assume that \(P(k)\) is true for some \(k\in\mathbb{Z}^+\).  
  Then for some integer \(m\),
  \[
  13^k+3\cdot 4^k+11=12m
  \]

  \textbf{Inductive step.}

  We must show that \(P(k+1)\) is true, i.e. that \(13^{k+1}+3\cdot 4^{k+1}+11\) is divisible by \(12\).

  Consider the difference:
  \begin{align*}
  &\left(13^{k+1}+3\cdot 4^{k+1}+11\right)-\left(13^k+3\cdot 4^k+11\right) \\
  &= 13^{k+1}-13^k+3\cdot 4^{k+1}-3\cdot 4^k \\
  &= 13^k(13-1)+3\cdot 4^k(4-1) \\
  &= 12\cdot 13^k+9\cdot 4^k \\
  &= 12\cdot 13^k+12\left(3\cdot 4^{k-1}\right)
  \end{align*}

  Since \(k\geq 1\), \(4^{k-1}\) is an integer, so this whole difference is a multiple of \(12\).

  Therefore
  \begin{align*}
  13^{k+1}+3\cdot 4^{k+1}+11
  &= \left(13^k+3\cdot 4^k+11\right)
   + \left[\left(13^{k+1}+3\cdot 4^{k+1}+11\right)-\left(13^k+3\cdot 4^k+11\right)\right] \\
  &= 12m + 12\cdot 13^k + 12\left(3\cdot 4^{k-1}\right) \\
  &= 12\left(m+13^k+3\cdot 4^{k-1}\right)
  \end{align*}

  So \(13^{k+1}+3\cdot 4^{k+1}+11\) is divisible by \(12\). Hence \(P(k+1)\) is true.

  Since \(P(1)\) is true and \(P(k)\Rightarrow P(k+1)\), it follows by mathematical induction that
  \[
  13^n+3\cdot 4^n+11
  \]
  is divisible by \(12\) for all \(n\in\mathbb{Z}^+\).

\end{enumerate}
\end{tcolorbox}


\newpage

% ---- Question 4 ----
\questionitem

Prove by mathematical induction that, for all non-negative integers $n$,
\[
2^n+2\cdot 3^{2n}+4^{2n+1}
\]
is divisible by $7$ \marks{6}

\vspace*{10pt}

\begin{tcolorbox}[colback=white, colframe=scarlet, title={\textbf{Solution}}, breakable, grow to left by=6mm, grow to right by=10mm]
We prove the statement by mathematical induction.

Let \(P(n)\) be the statement
\[
2^n + 2\times 3^{2n} + 4^{2n+1} \text{ is divisible by } 7
\]

\textbf{Base case: \(n=0\).}

\[
2^0 + 2\times 3^0 + 4^1 = 1+2+4=7
\]

Since \(7\) is divisible by \(7\), \(P(0)\) is true.

\textbf{Inductive hypothesis.}

Assume that \(P(k)\) is true for some non-negative integer \(k\).  
So
\[
2^k + 2\times 3^{2k} + 4^{2k+1}
\]
is divisible by \(7\). Hence, for some integer \(m\),
\[
2^k + 2\times 3^{2k} + 4^{2k+1} = 7m
\]

\textbf{Inductive step.}

We must show that \(P(k+1)\) is true, that is, that
\[
2^{k+1} + 2\times 3^{2k+2} + 4^{2k+3}
\]
is divisible by \(7\).

Starting with this expression,
\begin{align*}
2^{k+1} + 2\times 3^{2k+2} + 4^{2k+3}
&= 2\cdot 2^k + 2\cdot 3^2\cdot 3^{2k} + 4^2\cdot 4^{2k+1} \\
&= 2\cdot 2^k + 18\cdot 3^{2k} + 16\cdot 4^{2k+1} \\
&= 2\bigl(2^k + 2\times 3^{2k} + 4^{2k+1}\bigr) + 14\cdot 3^{2k} + 14\cdot 4^{2k+1}
\end{align*}

Now use the inductive hypothesis \(2^k + 2\times 3^{2k} + 4^{2k+1}=7m\):
\begin{align*}
2^{k+1} + 2\times 3^{2k+2} + 4^{2k+3}
&= 2(7m) + 14\cdot 3^{2k} + 14\cdot 4^{2k+1} \\
&= 14m + 14\bigl(3^{2k} + 4^{2k+1}\bigr) \\
&= 7\bigl(2m + 2\cdot 3^{2k} + 2\cdot 4^{2k+1}\bigr)
\end{align*}

This is \(7\) multiplied by an integer, so it is divisible by \(7\). Therefore \(P(k+1)\) is true.

Since \(P(0)\) is true and \(P(k)\Rightarrow P(k+1)\), it follows by mathematical induction that
\[
2^n + 2\times 3^{2n} + 4^{2n+1}
\]
is divisible by \(7\) for all non-negative integers \(n\).
\end{tcolorbox}


\newpage

% ---- Question 5 ----
\questionitem

For each positive integer $n$, define

\[u_n=5^{2n}-2^n\]

Use proof by induction to show that $23$ is a factor of $u_n$ for all $n \in \mathbb{Z}^+$ \marks{6}

\vspace*{10pt}

\begin{tcolorbox}[colback=white, colframe=scarlet, title={\textbf{Solution}}, breakable, grow to left by=6mm, grow to right by=10mm]
Let \(P(n)\) be the statement

\[
5^{2n}-2^n \text{ is divisible by } 23
\]

We prove \(P(n)\) for all \(n \in \mathbb{Z}^+\) by induction.

\textbf{Base case: \(n=1\).}

\[
u_1=5^{2}-2^1=25-2=23
\]

So \(u_1\) is divisible by \(23\). Hence \(P(1)\) is true.

\textbf{Inductive hypothesis.}

Assume that \(P(k)\) is true for some positive integer \(k\).  
So \(u_k\) is divisible by \(23\), which means there is an integer \(m\) such that

\[
u_k=5^{2k}-2^k=23m
\]

\textbf{Inductive step.}

Now consider \(u_{k+1}\):

\begin{align*}
u_{k+1}
&=5^{2(k+1)}-2^{k+1}\\
&=5^{2k+2}-2\cdot 2^k\\
&=25\cdot 5^{2k}-2\cdot 2^k\\
&=(23+2)5^{2k}-2\cdot 2^k\\
&=23\cdot 5^{2k}+2\cdot 5^{2k}-2\cdot 2^k\\
&=23\cdot 5^{2k}+2\left(5^{2k}-2^k\right)
\end{align*}

Using the inductive hypothesis \(5^{2k}-2^k=23m\),

\begin{align*}
u_{k+1}
&=23\cdot 5^{2k}+2(23m)\\
&=23\left(5^{2k}+2m\right)
\end{align*}

Since \(5^{2k}+2m\) is an integer, it follows that \(u_{k+1}\) is divisible by \(23\).  
So \(P(k+1)\) is true.

\textbf{Conclusion.}

Since \(P(1)\) is true, and \(P(k)\) true implies \(P(k+1)\) true, \(P(n)\) is true for all \(n\in\mathbb{Z}^+\) by mathematical induction.

Therefore, \(23\) is a factor of \(u_n=5^{2n}-2^n\) for all positive integers \(n\).
\end{tcolorbox}


\newpage

% ---- Question 6 ----
\questionitem

Prove by induction that, for $n \in \mathbb{Z}^+$,

\[
7^n-6n-1
\]

is divisible by $36$ \marks{6}

\vspace*{10pt}

\begin{tcolorbox}[colback=white, colframe=scarlet, title={\textbf{Solution}}, breakable, grow to left by=6mm, grow to right by=10mm]
Let \(P(n)\) be the statement
\[
7^n-6n-1 \text{ is divisible by } 36
\]

We prove \(P(n)\) is true for all \(n\in\mathbb{Z}^+\) by mathematical induction.

\textbf{Base case:} \(n=1\).

\[
7^1-6(1)-1=7-6-1=0
\]

Since \(0=36\times 0\), \(0\) is divisible by \(36\).  
So \(P(1)\) is true.

\textbf{Inductive hypothesis:} Assume that \(P(k)\) is true for some \(k\in\mathbb{Z}^+\).

So we assume
\[
7^k-6k-1
\]
is divisible by \(36\), which means there exists an integer \(m\) such that
\[
7^k-6k-1=36m
\]

\textbf{Inductive step:} We must show that \(P(k+1)\) is true, that is, that
\[
7^{k+1}-6(k+1)-1
\]
is divisible by \(36\).

Starting with this expression,
\begin{align*}
7^{k+1}-6(k+1)-1
&=7^{k+1}-6k-6-1\\
&=7^{k+1}-6k-7
\end{align*}

Now rewrite it to use the inductive hypothesis:
\begin{align*}
7^{k+1}-6k-7
&=7\cdot 7^k-6k-7\\
&=7(7^k-6k-1)+36k
\end{align*}

Using the inductive hypothesis \(7^k-6k-1=36m\),
\begin{align*}
7^{k+1}-6(k+1)-1
&=7(36m)+36k\\
&=36(7m+k)
\end{align*}

Since \(7m+k\in\mathbb{Z}\), this is a multiple of \(36\). Therefore
\[
7^{k+1}-6(k+1)-1
\]
is divisible by \(36\).

So \(P(k+1)\) is true.

Since \(P(1)\) is true, and \(P(k)\Rightarrow P(k+1)\), it follows by induction that
\[
7^n-6n-1 \text{ is divisible by } 36 \qquad \text{for all } n\in\mathbb{Z}^+
\]

Hence \(7^n-6n-1\) is divisible by \(36\) for all positive integers \(n\).
\end{tcolorbox}


\newpage

% ---- Question 7 ----
\questionitem

The function $f$ is defined by
\[
f(n)=9^n-8n+63 
\] 
Prove by induction that $f(n)$ is divisible by $64$ for $n \geq 1$ \marks{6}

\vspace*{10pt}

\begin{tcolorbox}[colback=white, colframe=scarlet, title={\textbf{Solution}}, breakable, grow to left by=6mm, grow to right by=10mm]
We prove by induction that
\[
9^n-8n+63 \text{ is divisible by } 64
\]
for all $n \ge 1$.

\textbf{Base case:} $n=1$.

\[
f(1)=9^1-8(1)+63=9-8+63=64
\]

Since $64$ is divisible by $64$, the statement is true for $n=1$.

\textbf{Induction hypothesis:}

Assume that for some integer $k \ge 1$,
\[
f(k)=9^k-8k+63
\]
is divisible by $64$.

So we may write
\[
9^k-8k+63=64m
\]
for some integer $m$.

\textbf{Induction step:}

Now consider $f(k+1)$:
\begin{align*}
f(k+1)
&=9^{k+1}-8(k+1)+63 \\
&=9\cdot 9^k-8k-8+63 \\
&=9\cdot 9^k-8k+55
\end{align*}

We now rewrite this in terms of $f(k)=9^k-8k+63$:
\begin{align*}
f(k+1)
&=9(9^k-8k+63)+64(k-8)
\end{align*}

To check this expansion:
\begin{align*}
9(9^k-8k+63)+64(k-8)
&=9^{k+1}-72k+567+64k-512 \\
&=9^{k+1}-8k+55 \\
&=f(k+1)
\end{align*}

Using the induction hypothesis $9^k-8k+63=64m$,
\begin{align*}
f(k+1)
&=9(64m)+64(k-8) \\
&=64(9m+k-8)
\end{align*}

Since $m$ and $k$ are integers, $9m+k-8$ is an integer, so $f(k+1)$ is divisible by $64$.

Therefore, if the statement is true for $n=k$, it is also true for $n=k+1$.

\textbf{Conclusion:}

The statement is true for $n=1$, and true for $n=k+1$ whenever it is true for $n=k$.

Hence, by mathematical induction,
\[
9^n-8n+63 \text{ is divisible by } 64
\quad \text{for all } n\ge 1
\]

So $f(n)$ is divisible by $64$ for all $n \ge 1$.
\end{tcolorbox}


\newpage

% ---- Question 8 ----
\questionitem

Prove by mathematical induction that
\[
7^{2n}+8\times 13^n-9
\]
is divisible by $36$ for all positive integers $n$ \marks{6}

\vspace*{10pt}

\begin{tcolorbox}[colback=white, colframe=scarlet, title={\textbf{Solution}}, breakable, grow to left by=6mm, grow to right by=10mm]
We prove the statement by mathematical induction.

Let
\[
E_n=7^{2n}+8\times 13^n-9
\]

We want to show that $E_n$ is divisible by $36$ for all positive integers $n$.

\textbf{Base case:} $n=1$.

\begin{align*}
E_1&=7^2+8\times 13-9\\
&=49+104-9\\
&=144\\
&=36\times 4
\end{align*}

So $E_1$ is divisible by $36$.

\textbf{Inductive hypothesis:} assume that for some positive integer $k$,
\[
E_k=7^{2k}+8\times 13^k-9
\]
is divisible by $36$.

So there exists an integer $q$ such that
\[
7^{2k}+8\times 13^k-9=36q
\]

\textbf{Inductive step:} consider $E_{k+1}$.

\begin{align*}
E_{k+1}
&=7^{2(k+1)}+8\times 13^{k+1}-9\\
&=7^{2k+2}+8\times 13^{k+1}-9\\
&=49\cdot 7^{2k}+104\cdot 13^k-9
\end{align*}

Now rewrite this so that $E_k$ appears:

\begin{align*}
E_{k+1}
&=13\cdot 7^{2k}+36\cdot 7^{2k}+13\cdot 8\cdot 13^k-117+108\\
&=13\left(7^{2k}+8\times 13^k-9\right)+36\left(7^{2k}+3\right)\\
&=13E_k+36\left(7^{2k}+3\right)
\end{align*}

Using the inductive hypothesis $E_k=36q$,

\begin{align*}
E_{k+1}
&=13(36q)+36\left(7^{2k}+3\right)\\
&=36\left(13q+7^{2k}+3\right)
\end{align*}

Since $13q+7^{2k}+3$ is an integer, $E_{k+1}$ is divisible by $36$.

So, if the statement is true for $n=k$, it is also true for $n=k+1$.

Since the statement is true for $n=1$, and true for $n=k$ implies true for $n=k+1$, it follows by mathematical induction that
\[
7^{2n}+8\times 13^n-9
\]
is divisible by $36$ for all positive integers $n$.
\end{tcolorbox}


\end{enumerate}
\end{document}